\section{The Lectionary cycles}
\label{sec:cycles}

\subsection{Two series, and their independence}

\begin{normtext}{\loc{olm} \latin{Praenotanda} 65}
Cursus lectionum in ``Proprio de tempore'' sic dispositus est: diebus dominicis et festis proponuntur
loci maioris momenti, ut partes praestantiores verbi Dei, intra congruum temporis spatium coram
fidelium coetu legi possint. Altera deinde series textuum sacrae Scripturae, quibus nuntius salutis
diebus festis explanatus quodammodo completur, proponitur pro feriis. Attamen harum praecipuarum
partium Ordinis lectionum neutra series, nempe dominicalis --- festiva et series ferialis, ex altera
pendet. Immo, Ordo lectionum dominicalis --- festivus per triennium explicatur, ferialis autem per
biennium. Unde Ordo lectionum dominicalis --- festivus seiunctus procedit a feriali, et vicissim.
\end{normtext}

\begin{workingrendering}
Neither of these principal parts of the Order of Readings --- the Sunday-and-festive series and the
ferial series --- depends on the other. Rather, the Sunday and festive Order of Readings runs over
three years, the ferial over two. Hence the Sunday and festive Order of Readings proceeds separately
from the ferial, and the ferial from it.
\end{workingrendering}

The independence is stated three times in one paragraph, in three constructions ---
\latin{neutra\dots ex altera pendet}, \latin{seiunctus procedit}, \latin{et vicissim} --- which is the
draftsman's way of saying that the point was expected to be missed. It is missed. The two commonest
errors in dated liturgical records are deriving a weekday numeral from a Sunday letter and carrying
a Sunday letter onto a weekday, and \loc{olm}~65 forecloses both. The periods are three and two: the
pairing of letter and numeral repeats only every six years, and in that interval every combination
occurs exactly once. Neither can ever determine the other.

\subsection{The Sunday cycle A, B, C}

\begin{normtext}{\loc{olm} \latin{Praenotanda} 66, footnote 102}
Singuli anni notantur litteris A, B, C. Ad statuendum autem quinam sit annus A, vel B, vel C, hoc
modo proceditur. Per litteram C designatur annus cuius numerus in tres partes aequales dividi
potest, quasi cyclus ab anno primo computi christiani initium sumpsisset. Sic annus 1 fuisset annus
A, annus 2 B, annus 3 annus C, et anni 6, 9, 12 \dots\ iterum annus C. Ita v.~g. annus 1980 est annus
C, annus vero sequens, scilicet annus 1981, est annus A, annus vero 1982 est annus B, et annus 1983
est iterum annus C. Et sic deinceps. Patet tamen singulos cyclos decurrere iuxta dispositionem anni
liturgici, nempe a prima hebdomada Adventus, quae cadit in anno civili praecedente.
\end{normtext}

The footnote gives a rule, a fiction to justify it, a worked series and a caveat, and all four are
needed.

The rule: the letter C designates a year whose number is divisible into three equal parts. The
fiction: as though the cycle had begun from the first year of the Christian era, so that year 1 would
have been A, year 2 B, year 3 C. The worked series: 1980 C, 1981 A, 1982 B, 1983 C. And the caveat,
already quoted in Section~\ref{sec:turn}: each cycle runs according to the arrangement of the
liturgical year, beginning from the First Week of Advent, which falls in the preceding civil year.

Writing $Y$ for the civil year that the footnote numbers --- the year in which the liturgical year
\emph{ends}, that is the year of its January to November, not the year of its Advent:

\begin{itemize}
\item remainder 1 on division of $Y$ by 3 gives Year A;
\item remainder 2 gives Year B;
\item remainder 0 --- that is, exact divisibility --- gives Year C.
\end{itemize}

Equivalently, the cycle beginning at the First Sunday of Advent in civil year $y$ is fixed by the
remainder of $y+1$ on division by 3. The footnote's own series satisfies this: 1980 divides exactly,
C; 1981 leaves 1, A; 1982 leaves 2, B; 1983 divides exactly, C.

Two properties of the rule are worth naming because they are where errors enter. It is a property of
the \emph{civil} year number, but it labels a \emph{liturgical} year that straddles two civil years;
so one civil year always carries the tail of one cycle and the head of the next, and any record that
states a letter without stating an interval is ambiguous. And it turns at the First Sunday of Advent
and at no other moment: not at 1 January, not at Easter, not at the beginning of Ordinary Time.

The Missal's own \latin{Tabella temporaria} confirms the arithmetic independently. Its \latin{Cycli
domin.} column prints a pair for each civil year --- 2001 as \latin{C - A}, 2002 as \latin{A - B},
2003 as \latin{B - C}, and so on --- which is precisely the tail-and-head structure the footnote
implies. All twenty-four rows agree with the formula above.

The second paragraph of the footnote adds the association that gives the cycle its shape: the years
take their character from the synoptic gospel read semicontinuously in Ordinary Time, so that the
first year of the cycle is the year of reading Matthew, the second of reading Mark and the third of
reading Luke.

\subsection{The weekday cycle I and II}

\begin{normtext}{\loc{olm} \latin{Praenotanda} 69, 4}
Pro feriis vero triginta quattuor hebdomadarum ``per annum'' lectiones evangelicae unico disponuntur
cyclo, qui singulis annis resumitur. Prior vero lectio, in duplici cyclo ordinatur, alternis annis
legenda. Annus primus adhibetur annis imparibus, annus secundus vero annis paribus.
\end{normtext}

\begin{workingrendering}
For the weekdays of the thirty-four weeks of Ordinary Time the gospel readings are arranged in a
single cycle, repeated each year. The first reading, however, is arranged in a double cycle, to be
read in alternate years. Year one is used in odd-numbered years, year two in even-numbered years.
\end{workingrendering}

Three things are settled here that are commonly stated loosely.

First, only the first reading alternates. The weekday gospel of Ordinary Time is a single cycle
repeated every year; Cycle~I and Cycle~II differ in the first reading alone. A record that describes
``the Cycle II gospel'' has misdescribed its own source.

Second, the parity is of the same $Y$ as the Sunday cycle --- the year in which the liturgical year
ends --- because both cycles turn together at the First Sunday of Advent. Odd is Cycle~I, even is
Cycle~II. The turn is simultaneous; the values are unrelated.

Third, the article speaks of the weekdays \emph{of the thirty-four weeks of Ordinary Time} and of
nothing else. The rest of \loc{olm}~69 disposes of the other seasons: Lent has a one-year cycle
ordered on its own principles, baptismal and penitential (69.2); the weekdays of Advent, Christmas
Time and Easter Time are likewise annual, and their readings do not change (69.3). So the numeral
I or II is meaningless outside Ordinary Time, and a dated record for a weekday of Lent or Easter that
carries a weekday-cycle numeral has recorded a category error, not a fact.

\subsection{The promulgated starting point}

Both formulas can be checked at their origin, because the decree that put the Order of Readings into
force named the first year of each cycle. The decree of the Sacred Congregation for Divine Worship,
protocol 106/69, dated 25 May 1969, Pentecost Sunday, is reprinted in the front matter of the 1981
edition:

\begin{normtext}{\loc{olm}, \latin{Decretum} of 25 May 1969, protocol 106/69}
Sacra igitur haec Congregatio pro Cultu Divino, de speciali mandato Summi Pontificis, eundem Ordinem
lectionum Missae promulgat, statuens ut a die 30 novembris, dominica prima Adventus, huius anni 1969
vigere incipiat. Adhibebitur autem proximo anno liturgico, series B lectionum dominicalium et series
II prioris lectionis in feriis ``per annum''.
\end{normtext}

\begin{workingrendering}
This Sacred Congregation for Divine Worship therefore, by special mandate of the Supreme Pontiff,
promulgates the same Order of Readings for Mass, establishing that it begin to have force from 30
November, the First Sunday of Advent, of this year 1969. In the coming liturgical year, series B of
the Sunday readings and series II of the first reading on the weekdays of Ordinary Time will be used.
\end{workingrendering}

The liturgical year beginning on 30 November 1969 ends in civil 1970. Then $Y = 1970$;
1970 leaves remainder 2 on division by 3, which is Year~B; and 1970 is even, which is Cycle~II. Both formulas reproduce the
promulgated starting point exactly, from a decree written twelve years before the footnote that
states the rule. The two cycles have run without interruption from that Sunday, and every letter and
numeral in this reference is, in the end, a count forward from it.

\subsection{Where neither cycle governs}

The remaining parts of the Order of Readings run by their own laws, as \loc{olm}~65 says in its last
paragraph and \loc{olm}~70--72 spell out. Celebrations of saints have a double series: proper
readings for solemnities, feasts and memorials where they exist, and the far larger provision in the
Commons, from which a choice is made unless otherwise expressly noted (70--71). Ritual Masses,
Masses for various needs, votive Masses and Masses for the dead have the same arrangement --- several
texts proposed together, as in the Commons (72). None of these is cycle-governed, and a record for
such a celebration should read \emph{not applicable} rather than carry a letter or a numeral.

\subsection{Reading a cycle statement correctly}

The Secretariat of Divine Worship's annual calendar states both cycles for each year, and states
them in different units, which is itself instructive. For the liturgical year 2026 it gives the
Sunday cycle as Year~A from 30~November 2025 to 22~November 2026, and the weekday cycle as Cycle~2
from 12~January to 17~February 2026 and again from 25~May to 28~November 2026.

The Sunday interval is bounded by Sundays: it ends on Christ the King, the last Sunday the cycle
governs. The weekday intervals are bounded by weekdays and are given as two runs, because the ferial
Ordinary Time cycle is interrupted by Lent and Easter Time and does not govern them: it begins on
the Monday after the Baptism of the Lord, stops on the Tuesday before Ash Wednesday, and resumes on
the Monday after Pentecost, running to the Saturday before Advent. Neither statement is the
liturgical year's own boundary, and neither is wrong. Both describe the extent over which the cycle
they name actually operates.

Three rules of hygiene follow, and they are the ones this reference has found most often broken in
secondary sources:

\begin{itemize}
\item record a cycle letter with the exact interval it governs, never bare;
\item resolve the weekday numeral independently of the Sunday letter, from the parity of $Y$; and
\item assign no cycle at all to a celebration whose readings are not cycle-governed.
\end{itemize}
